\subsection{DEFINITION OF A ROOT SYSTEM}

Lemma 1. Let V be a vector space over k, R a finite subset of V generating V. For any \(\alpha \in R\) such that \(\alpha \neq 0\) , there exists at most one reflection s of V such that \(s(\alpha) = -\alpha\) and \(s(R) = R\) .

Let G be the group of automorphisms of V leaving R stable. Since R generates V, G is isomorphic to a subgroup of the symmetric group of R, and hence is finite. Let \(s, s'\) be reflections of V such that \(s(\alpha) = s'(\alpha) = -\alpha\) , \(s(R) = R\) , \(s'(R) = R\) . Then \(t = ss'\) belongs to G, and hence is of finite order m. On the other hand,

\[
t (\alpha) = \alpha \text { and } t (x) \equiv x \bmod k \alpha \text { for   all } x \in V.
\]

Hence, there exists a linear form \(f\) on \(V\) such that

\[
t (x) = x + f (x) \alpha \text {   for   all   } x \in V
\]

and \(f(\alpha) = 0\) . By induction on \(n\) , it follows that

\[
t ^ {n} (x) = x + n f (x) \alpha \text {   for   all   } x \in V.
\]

Taking \(n\) equal to \(m\) , we see that \(mf(x) = 0\) for all \(x \in V\) , so \(f = 0\) , \(t = 1\) and \(s = s'\) .

DEFINITION 1. Let V be a vector space over k, and R a subset of V. Then R is said to be a root system in V if the following conditions are satisfied:

(RS \(_{I}\) ) R is finite, does not contain 0, and generates V.

(RS \(_{II}\) ) For all \(\alpha \in \mathbb{R}\) , there exists an element \(\alpha^{\prime}\) of the dual V \(^{*}\) of V such that \(\langle \alpha, \alpha^{\prime} \rangle = 2\) and that the reflection \(s_{\alpha, \alpha'}\) (cf.~Chap. V, § 2) leaves R stable.

(RS \(_{III}\) ) For all \(\alpha \in R\) , \(\alpha^{\sim}(R) \subseteq \mathbf{Z}\) .

By Lemma 1, the reflection \(s_{\alpha, \alpha'}\) (and hence also the linear form \(\alpha'\) ) is determined uniquely by \(\alpha\) , so \((\mathrm{RS}_{\mathrm{III}})\) makes sense. We put \(s_{\alpha, \alpha'} = s_{\alpha}\) . Then \(s_{\alpha}(x) = x - \langle \alpha', x \rangle \alpha\) for all \(x \in V\) .

The elements of R are called the roots (of the system considered). The dimension of V is called the rank of the system.

The automorphisms of V that leave R stable are called the automorphisms of R. They form a finite group denoted by A(R). The subgroup of A(R) generated by the \(s_{\alpha}\) is called the Weyl group of R and is denoted by W(R), or simply by W.

Remark. 1) Let \(k'\) be an extension of \(k\) . Identify V canonically with a subset of \(V \otimes k'\) and \(V^*\) with a subset of \(V^* \otimes k' = (V \otimes k')^*\) . Then, R is a root system in \(V \otimes k'\) , and the \(\alpha'\) are the same as before.

Lemma 2. Let R be a root system in V. Let \((x|y)\) be a symmetric bilinear form on V, non-degenerate and invariant under W(R). Identify V with \(V^{*}\) by means of this form. If \(\alpha \in R\) , then \(\alpha\) is non-isotropic and

\[
\alpha^ {\prime} = \frac {2 \alpha}{(\alpha | \alpha)}.
\]

This follows from formula (4) of Chap. V, § 2, no. 3.

PROPOSITION 1. Let \(\mathrm{V}_{\mathbf{Q}}\) (resp. \(\mathrm{V}_{\mathbf{Q}}^{*}\) ) be the \(\mathbf{Q}\) -vector subspace of \(\mathrm{V}\) (resp. \(\mathrm{V}^{*}\) ) generated by the \(\alpha\) (resp. the \(\alpha^{-}\) ). Then \(\mathrm{V}_{\mathbf{Q}}\) (resp. \(\mathrm{V}_{\mathbf{Q}}^{*}\) ) is a \(\mathbf{Q}\) -structure on \(\mathrm{V}\) (resp. \(\mathrm{V}^{*}\) ) (Algebra, Chap. II, § 8, no. 1). The restriction to \(\mathrm{V}_{\mathbf{Q}} \times \mathrm{V}_{\mathbf{Q}}^{*}\) of the canonical bilinear form on \(\mathrm{V} \times \mathrm{V}^{*}\) gives an identification of each of the spaces \(\mathrm{V}_{\mathbf{Q}}, \mathrm{V}_{\mathbf{Q}}^{*}\) with the dual of the other. The set \(\mathbb{R}\) is a root system in \(\mathrm{V}_{\mathbf{Q}}\) .

If \(k = \mathbf{R}\) , there exists a scalar product on \(V\) invariant under \(W(R)\) (Integration, Chap. VII, § 3, no. 1, Prop. 1); Lemma 2 now shows that the \(\alpha^{-}\) generate \(V^{*}\) . By Remark 1, the \(\alpha^{-}\) again generate \(V^{*}\) if \(k = \mathbf{Q}\) . We now go to the general case. Put \(E = V_{\mathbf{Q}}\) . By \((RS_{III})\) , each \(\alpha^{-}\) maps \(E\) to \(\mathbf{Q}\) , and so defines an element \(\tilde{\alpha}\) of \(E^{*}\) . It is immediate that \(R\) is a root system in \(E\) , and that the element corresponding to \(\alpha\) in \(E^{*}\) is \(\tilde{\alpha}\) . By what we said above, the \(\tilde{\alpha}\) generate the vector space \(E^{*}\) . Consider the canonical homomorphism \(i: E \otimes_{\mathbf{Q}} k \to V\) , and its transpose \({}^{t}i: V^{*} \to E^{*} \otimes_{\mathbf{Q}} k\) . Since \(R\) generates \(V\) , \({}^{t}i\) is injective; but the image of \({}^{t}i\) contains the \(\tilde{\alpha}\) , so \({}^{t}i\) is surjective. From this we conclude finally that \(i\) and \({}^{t}i\) are isomorphisms. We can therefore identify \(V\) with \(E \otimes k\) , \(V^{*}\) with \(E^{*} \otimes k\) , \(\alpha^{-}\) with \(\tilde{\alpha}\) , and \(V_{\mathbf{Q}}^{*}\) with \(E^{*}\) . Thus, \(V_{\mathbf{Q}}\) (resp. \(V_{\mathbf{Q}}^{*}\) ) is a \(Q\) -structure on \(V\) (resp. \(V^{*}\) ). The restriction to \(V_{\mathbf{Q}} \times V_{\mathbf{Q}}^{*}\) of the canonical bilinear form on \(V \times V^{*}\) can be identified with the canonical bilinear form on \(E \times E^{*}\) , hence the proposition.

Remarks. 2) Proposition 1 reduces the study of root systems to the case \(k = \mathbf{Q}\) . Remark 1 reduces it further to the study of root systems in the real vector space \(V_{\mathbf{R}} = V_{\mathbf{Q}} \otimes_{\mathbf{Q}} \mathbf{R}\) . The Weyl groups associated to these different systems are canonically identified.

\begin{enumerate}
\def\labelenumi{\arabic{enumi})}
\setcounter{enumi}{2}
\tightlist
\item
  Since the \(\alpha^{\sim}\) generate \(V^{*}\) , the group \(W(R)\) , considered as a subgroup of \(\mathbf{GL}(V_{\mathbf{R}})\) , is essential (Chap. V, § 3, no. 7). Moreover, the Cor. of Th. 1 of Chap. V, § 3, no. 2 shows that the only reflections belonging to \(W(R)\) are the \(s_{\alpha}\) .
\end{enumerate}

PROPOSITION 2. The \(\alpha^{\vee}\) form a root system in \(V^{*}\) , and \(\alpha^{\vee} = \alpha\) for all \(\alpha \in R\) .

The \(\alpha^{\prime}\) satisfy \((\mathrm{RS}_{\mathrm{I}})\) by Prop. 1. Since \(s_{\alpha, \alpha'}\) is an automorphism of the vector space V equipped with the subset R, \({}^{t}(s_{\alpha, \alpha'})^{-1}\) leaves the set R' of the \(\alpha'\) stable; but \({}^{t}(s_{\alpha, \alpha'})^{-1} = s_{\alpha', \alpha}\) , which proves that R' satisfies \((\mathrm{RS}_{\mathrm{II}})\) and that \(\alpha' = \alpha\) . Finally, \(\langle \alpha', \beta \rangle \in \mathbf{Z}\) for all \(\alpha' \in \mathrm{R'}\) and \(\beta \in \mathrm{R}\) , so R' satisfies \((\mathrm{RS}_{\mathrm{III}})\) .

The set \(R^{-}\) is called the inverse root system of R. The map \(\alpha \mapsto \alpha^{-}\) is a bijection from R to \(R^{-}\) , called the canonical bijection from R to \(R^{-}\) . Note that, if \(\alpha, \beta\) are elements of R such that \(\alpha + \beta \in R\) , then \((\alpha + \beta)^{-} \neq \alpha^{-} + \beta^{-}\) in general.

Since \(s_{\alpha}(\alpha) = -\alpha\) , axiom (RSII) shows that \(-R = R\) . Evidently \((- \alpha)^{\sim} = -\alpha^{\sim}\) and \(-1 \in A(R)\) (but it is not always true that \(-1 \in W(R)\) ).

The equality \({}^{t}(s_{\alpha,\alpha^{-}})^{-1}=s_{\alpha^{-},\alpha}\) shows that the map \(u\mapsto{}^{t}u^{-1}\) is an isomorphism from the group W(R) to the group W(R \(^{*}\) ). We identify these two groups by means of this isomorphism; in other words, we consider W(R) as acting both in V and in V \(^{*}\) . Similarly for A(R).

PROPOSITION 3. For \(x, y \in V\) , put

\[
(x | y) = \sum_ {\alpha \in R} \langle \alpha^ {\vee}, x \rangle \langle \alpha^ {\vee}, y \rangle .
\]

Then \((x|y)\) is a non-degenerate symmetric bilinear form on \(\mathbf{V}\) , invariant under \(\mathbf{A}(\mathbb{R})\) . For \(x, y \in \mathrm{V}_{\mathbf{Q}}\) we have \((x|y) \in \mathbf{Q}\) . The canonical extension of \((x|y)\) to

\[
V _ {R} = V _ {Q} \otimes_ {Q} R
\]

is non-degenerate and positive.

It is clear that \((x|y)\) is a symmetric bilinear form on V. If \(g\in \mathrm{A}(\mathbb{R})\)

\[
(g (x) | g (y)) = \sum_ {\alpha \in R} \left\langle {} ^ {t} g \left(\alpha^ {-}\right), x \right\rangle \left\langle {} ^ {t} g \left(\alpha^ {-}\right), y \right\rangle = (x | y)
\]

since \((^t g)(\mathbb{R}^\sim) = \mathbb{R}^\sim\) . If \(x, y \in \mathrm{V}_{\mathbf{Q}}\) , then \((x|y) \in \mathbf{Q}\) by \((\mathrm{RS}_{\mathrm{III}})\) . If \(z \in \mathrm{V}_{\mathbf{R}}\) , then \((z|z) = \sum_{\alpha \in \mathbb{R}} \langle \alpha^\sim, z \rangle^2 \geqslant 0\) , and \((z|z) > 0\) if \(z \neq 0\) by Prop. 1, so the canonical extension of \((x|y)\) to \(V_{R}\) is positive and non-degenerate. The restriction of \((x|y)\) to \(V_{Q}\) is thus non-degenerate, and hence the form \((x|y)\) on V is non-degenerate.

PROPOSITION 4. (i) Let X be a subset of R, let \(V_{X}\) be the vector subspace of V generated by X, and let \(V'_{X}\) be the vector subspace of \(V^{*}\) generated by the \(\alpha'\) , where \(\alpha \in X\) . Then V is the direct sum of \(V_{X}\) and the orthogonal complement of \(V'_{X}\) , \(V^{*}\) is the direct sum of \(V'_{X}\) and the orthogonal complement of \(V_{X}\) , and \(V'_{X}\) is identified with the dual of \(V_{X}\) .

\begin{enumerate}
\def\labelenumi{(\roman{enumi})}
\setcounter{enumi}{1}
\tightlist
\item
  \(R \cap V_{X}\) is a root system in \(V_{X}\) , and the canonical bijection from \(R \cap V_{X}\) to its inverse root system is identified with the restriction of the map \(\alpha \mapsto \alpha^{-}\) to \(R \cap V_{X}\) .
\end{enumerate}

By Remark 2, we can assume that \(k = \mathbf{R}\) . Identify \(V\) with \(V^*\) by means of the symmetric bilinear form of Prop. 3. We have \(\alpha^{-} = \frac{2\alpha}{(\alpha|\alpha)}\) for all \(\alpha \in R\) (Lemma 2). Every vector subspace of \(V\) is non-isotropic, and the proposition is now clear.

COROLLARY. Let \(V_{1}\) be a vector subspace of V, and let \(V_{2}\) be the vector subspace generated by \(R \cap V_{1}\) . Then \(R \cap V_{1}\) is a root system in \(V_{2}\) .

This follows from (ii) applied to \(\mathrm{X} = \mathrm{R} \cap \mathrm{V}_1\) .

For \(\alpha, \beta \in \mathbb{R}\) , put

\[
\langle \alpha , \beta^ {\prime} \rangle = n (\alpha , \beta).\tag{1}
\]

Then

\[
n (\alpha , \alpha) = 2\tag{2}
\]

\[
n (- \alpha , \beta) = n (\alpha , - \beta) = - n (\alpha , \beta).\tag{3}
\]

By \((\mathrm{RS}_{\mathrm{III}})\)

\[
n (\alpha , \beta) \in \mathbf {Z}.\tag{4}
\]

By the definition of \(n(\alpha, \beta)\) ,

\[
s _ {\beta} (\alpha) = \alpha - n (\alpha , \beta) \beta .\tag{5}
\]

Formula (1) and Prop. 2 imply that

\[
n (\alpha , \beta) = n \left(\beta^ {-}, \alpha^ {-}\right).\tag{6}
\]

Let \((x|y)\) be a symmetric bilinear form on V, non-degenerate and invariant under W(R) (Prop. 3). By Lemma 2,

\[
n (\alpha , \beta) = \frac {2 (\alpha | \beta)}{(\beta | \beta)}.\tag{7}
\]

It follows that

\[
n (\alpha , \beta) = 0 \Longleftrightarrow n (\beta , \alpha) = 0 \Longleftrightarrow (\alpha | \beta) = 0 \Longleftrightarrow s _ {\alpha} \text {   and   } s _ {\beta} \text {   commute.   (8) }
\]

\[
\text {   If   } (\alpha | \beta) \neq 0, \text {   then   } \frac {n (\beta , \alpha)}{n (\alpha , \beta)} = \frac {(\beta | \beta)}{(\alpha | \alpha)}.\tag{9}
\]

